\documentclass[11pt]{article}

\usepackage[margin=1in]{geometry}
\usepackage{amsmath,amssymb,amsthm}
\usepackage{mathtools}
\usepackage{hyperref}
\usepackage{xcolor}
\usepackage{enumitem}

\hypersetup{colorlinks=true, linkcolor=blue, citecolor=blue, urlcolor=blue}

\newcommand{\R}{\mathbb{R}}
\newcommand{\N}{\mathbb{N}}
\newcommand{\Rstar}{\mathbb{R}^{\star}}
\newcommand{\eps}{\varepsilon}
\let\straightepsilon\epsilon
\renewcommand{\epsilon}{\varepsilon}
\newcommand{\st}{\operatorname{st}}
\newcommand{\Var}{\operatorname{Var}}

\setlist[description]{leftmargin=0pt, style=unboxed, font=\normalfont\bfseries}

\title{Twenty Exercises in Algebraic Probability\\
\large Algebraic Stochastics and Statistics over $\Rstar$}
\author{Exercises for engineers, with a machine-checked Lean~4 companion}
\date{\today}

\begin{document}
\maketitle


\section*{The one rule you need}

A probability is an ordinary algebraic value in the hyperreal field $\Rstar$,
not a real number obtained by a limit. Fix a positive infinitesimal $\eps$ and
put $\omega = 1/\eps$, so that
\[
  \eps \cdot \omega = 1,
  \qquad 0 < \eps < r \quad\text{for every real } r > 0 .
\]
On a uniform hyperfinite sample space $\Omega$ the probability of an event is a
plain counting ratio,
\[
  P(A) = \frac{\#A}{\#\Omega},
\]
so if the experiment has $A\omega^{n}$ outcomes and the event has $c\omega^{d}$
favorable ones, then
\[
  P(E) = \frac{c\,\omega^{d}}{A\,\omega^{n}} = \frac{c}{A}\,\eps^{\,n-d}.
\]
Everything else is built from these ratios by addition, multiplication,
complement and division. The standard part $\st$ is applied only when the
ordinary real shadow of an answer is explicitly wanted --- which is exactly the
step that classically destroys the information these exercises are about.

\paragraph{Theory.} The axioms, the model, and the underlying calculus are
\emph{not} repeated here. They are developed in the companion paper,
\href{https://files.pannous.com/hyperreals.pdf}{\emph{Hyperreal Numbers: $\eps\cdot\omega=1$ --- An
Algebraic, Computable Introduction to Infinitesimal Calculus}}
(\url{https://files.pannous.com/hyperreals.pdf}), with the machine-checked Lean~4 development at
\url{https://github.com/pannous/hyper-lean}.

\paragraph{Readiness labels.} Each exercise is tagged by what the present
Lean~4 formalization already supports.
\begin{description}
\item[Now] reduces to arithmetic and order facts already available for the
concrete $\Rstar$ representation, even if a polished event API is still absent.
\item[Partial] the scalar answer is representable now, but reusable random
variables, hyperfinite counts, or finite-sum infrastructure are missing.
\item[Research] deliberately requires a substantive extension of the framework.
\end{description}


\section*{Exercises}


\subsection*{Exercise 1 --- One exact outcome \hfill \normalfont\small[Now]}

\textbf{Problem.} Let $\Omega = \{0, 1, ..., \omega - 1\}$ be uniform, and let $j$ be
one specified outcome. Compute $P(\{j\})$, prove that it is positive, and compute
its standard part.

\textbf{Why this framework.} An ordinary continuum model records only the shadow
$P(\{j\}) = 0$. Hyperfinite counting records simultaneously that the outcome is
possible and that its ordinary shadow is zero.


\subsection*{Exercise 2 --- A finite target \hfill \normalfont\small[Now]}

\textbf{Problem.} In the same $\Omega$, let $A$ contain exactly $r$ specified
outcomes, where $r$ is a positive standard natural number. Find $P(A)$ and
compare it with a singleton probability.

\textbf{Why this framework.} The ordinary shadow assigns zero to both events. The
algebraic answer retains the finite likelihood ratio between them.


\subsection*{Exercise 3 --- Union and overlap without collapsing to zero \hfill \normalfont\small[Now/Partial]}

\textbf{Problem.} Suppose finite events $A,B\subset \Omega$ satisfy $\#A = a$,
$\#B = b$, and $\#(A\cap B) = c$, with standard finite $a,b,c$. Compute
the probability of their union.

\textbf{Why this framework.} Taking standard parts first turns all four small-event
probabilities into zero, hiding the overlap correction. Algebra in $\Rstar$ keeps
ordinary inclusion-exclusion informative at infinitesimal scale.


\subsection*{Exercise 4 --- The dart: point versus segment \hfill \normalfont\small[Now]}

\textbf{Problem.} Discretize the unit square by an $\omega$ by $\omega$ uniform grid.
A specified point occupies one grid site. A horizontal segment of standard
length $L > 0$ contains $L\cdot \omega$ sites (choose a compatible hyperfinite grid).
Compute both probabilities and their ratio.

\textbf{Why this framework.} In the ordinary square both events have probability
zero, so their very different rarity orders become indistinguishable. The
algebraic grid explains the distinction by counting, without introducing a
separate set-valued foundation.


\subsection*{Exercise 5 --- Coefficients cannot defeat a rarity order \hfill \normalfont\small[Now]}

\textbf{Problem.} On the same square, let $A$ consist of $M = 10^{100}$ specified
points and let $B$ be one complete row, disjoint from $A$. Compare $P(A)$ and
$P(B)$, and compute $P(A\cup B)$.

\textbf{Why this framework.} Both events again have ordinary shadow zero. Merely
saying that $A$ has an enormous finite number of outcomes misses that one full
row belongs to a strictly larger infinitesimal order.


\subsection*{Exercise 6 --- Near certainty with visible failure \hfill \normalfont\small[Now]}

\textbf{Problem.} From the one-dimensional $\Omega$, choose uniformly after marking
$r$ excluded outcomes. Find the probability $G$ of choosing an allowed
outcome and the probability of failure.

\textbf{Why this framework.} The ordinary shadow reports $P(G)=1$ and failure
probability zero. The algebraic value distinguishes certainty from an event
that can fail in exactly $r$ exceptional ways.


\subsection*{Exercise 7 --- Conditioning on a finite rare event \hfill \normalfont\small[Now]}

\textbf{Problem.} Let $C$ consist of $k > 0$ specified outcomes of the
one-dimensional $\Omega$, and let $j\in C$. Compute $P(\{j\}\mid C)$ using the
algebraic definition $P(A\mid B) = P(A\cap B)/P(B)$.

\textbf{Why this framework.} After taking ordinary shadows, numerator and
denominator are both zero and the elementary ratio is unavailable. The
infinitesimal factors cancel before any information is discarded.


\subsection*{Exercise 8 --- Conditioning a point on a line \hfill \normalfont\small[Now]}

\textbf{Problem.} In the $\omega$ by $\omega$ square, let $R$ be a specified complete
row and let $p$ be one specified point on it. Find $P(\{p\}\mid R)$.

\textbf{Why this framework.} The ordinary two-dimensional shadows of both $R$ and
$\{p\}$ are zero, so a direct real-valued quotient reads $0/0$. Algebraic
probability retains their codimension difference and cancels it correctly.


\subsection*{Exercise 9 --- Independence of coordinate events \hfill \normalfont\small[Now/Partial]}

\textbf{Problem.} Choose $(X,Y)$ uniformly from the $\omega$ by $\omega$ grid. Let
$A = \{X=x0\}$ and $B = \{Y=y0\}$. Verify independence algebraically.

\textbf{Why this framework.} Ordinary shadows turn $P(A)$, $P(B)$, and
$P(A\cap B)$ into zero, so the equation $0 = 0\cdot 0$ cannot explain the
finite-grid structure. The infinitesimal equality is nontrivial.


\subsection*{Exercise 10 --- Dependence among equally rare events \hfill \normalfont\small[Now]}

\textbf{Problem.} Work on the toroidal $\omega$ by $\omega$ grid. Let
$D_{0} = \{Y=X\}$ and $D_{1} = \{Y=X+1\}$, and put $A=D_{0}$, $B=D_{0}\cup D_{1}$. Are $A$ and
$B$ independent? Compute $P(A\mid B)$.

\textbf{Why this framework.} All of $A$, $B$, and their intersection have ordinary
shadow zero. Those shadows erase both the strong dependence and the exact
conditional probability.


\subsection*{Exercise 11 --- Moments of a rare Bernoulli variable \hfill \normalfont\small[Now/Partial]}

\textbf{Problem.} Let $X$ equal $1$ with probability $q=c\cdot \eps$ and $0$
otherwise, for a positive standard $c$. Compute $E[X]$ and $\Var(X)$ exactly.

\textbf{Why this framework.} Replacing $q$ by its ordinary shadow makes $X$
identically zero and deletes its first- and second-order behavior. Algebraic
moments retain both.


\subsection*{Exercise 12 --- An infinitesimal chance with finite expectation \hfill \normalfont\small[Now/Partial]}

\textbf{Problem.} A payoff $Y$ equals $\omega/c$ when the rare event from Exercise 11
occurs and equals zero otherwise. Compute its expectation and variance.

\textbf{Why this framework.} Sending the event probability to zero before
multiplying by the infinite payoff gives an indeterminate or misleading
picture. The gauged algebra performs the cancellation exactly and also shows
that finite expectation need not imply finite variance.


\subsection*{Exercise 13 --- Exact finite-resolution corrections to uniform moments \hfill \normalfont\small[Partial]}

\textbf{Problem.} Let $J$ be uniform on $\{0,...,\omega-1\}$ and set
$X=J/\omega$. Compute $E[X]$, $E[X^2]$, and $\Var(X)$ without taking a limit.

\textbf{Why this framework.} The ordinary uniform law gives only $1/2$ and $1/12$.
The algebraic result also records how the chosen grid convention approaches
those values.


\subsection*{Exercise 14 --- Collision statistic \hfill \normalfont\small[Partial]}

\textbf{Problem.} Draw $n$ independent labels uniformly from an $\omega$-element
set, where $n$ is standard and finite. Let $C$ count equal unordered pairs.
Find $E[C]$ and $\Var(C)$.

\textbf{Why this framework.} A continuous classical shadow says that ties never
occur and makes both answers zero. The algebraic answer retains the leading
collision rate and its correction, which matters when comparing rare data
quality failures.


\subsection*{Exercise 15 --- An estimator with infinitesimal contamination \hfill \normalfont\small[Partial]}

\textbf{Problem.} Let $X_{1},...,X_{n}$ be independent Bernoulli variables with
$P(X_{i}=1)=q=\theta+a\cdot \eps$, where $0<\theta<1$, $a$ is standard, and $n$ is
standard and positive. For the sample mean $T$, compute its expectation and
variance through order $\eps^2$.

\textbf{Why this framework.} The ordinary shadow cannot distinguish contamination
$+a\cdot \eps$ from none at all. Algebraic statistics retains its direction and
its effect on both bias and uncertainty.


\subsection*{Exercise 16 --- An infinitesimal likelihood comparison \hfill \normalfont\small[Now/Partial]}

\textbf{Problem.} A Bernoulli sample has $s$ successes and $f$ failures, both
standard finite. Compare the candidate parameters $q$ and $q+\eps$, where
$0<q<1$, by their exact likelihood ratio. Which candidate wins when the first
nonzero infinitesimal coefficient is of order $\eps$?

\textbf{Why this framework.} Their real shadows are the same parameter $q$, so a
classical parameter representation declares an immediate tie. The algebraic
likelihood orders candidates that differ below ordinary resolution.


\subsection*{Exercise 17 --- Bayes conditioning when every observed route is rare \hfill \normalfont\small[Research]}

\textbf{Problem.} A condition $D$ has probability $\eps$. A test is always
positive under $D$ and has false-positive probability $c\cdot \eps$ under
$D^c$, for standard $c>0$. Compute $P(D\mid +)$ exactly and find its
standard part.

\textbf{Why this framework.} If both prevalence and false-positive probability are
replaced by zero first, the elementary Bayes quotient loses their relative
scale. Keeping them algebraic produces a finite, informative posterior from
two rare routes to the observation.


\subsection*{Exercise 18 --- An extreme exact test with $\omega$ trials \hfill \normalfont\small[Research]}

\textbf{Problem.} Under a fair Bernoulli null hypothesis, conduct $\omega$ trials
and observe success every time. Give the exact one-sided probability of an
outcome at least this extreme and compare it with $\eps^k$ for every
standard natural $k$.

\textbf{Why this framework.} The usual infinite-trial limit reports zero. An
algebraic extension should retain a strictly positive value and show that
exponential rarity is smaller than every fixed polynomial order of
$\eps$.


\subsection*{Exercise 19 --- Poisson probabilities from one hyperfinite experiment \hfill \normalfont\small[Research]}

\textbf{Problem.} Perform $\omega$ independent Bernoulli trials with success
probability $\lambda\cdot \eps$, where $\lambda>0$ is standard. For fixed standard
$k$, compute $P(K=k)$ algebraically and identify its standard part.

\textbf{Why this framework.} Classical derivations introduce a sequence of
binomial experiments and then take a limit. The hyperfinite formulation puts
the large count and tiny probability into one experiment, with
$\omega\cdot \eps=1$, and retains corrections beyond the limiting law.


\subsection*{Exercise 20 --- Uniform local asymptotic statistics in one algebraic scale \hfill \normalfont\small[Research]}

\textbf{Problem.} Fix standard $p\in (0,1)$ and put
$v=p(1-p)$, $\sigma=\sqrt{v}$, $N=\omega$, and $\eps=N^{-1}$. For standard
$h$, let $p_h=p+h\cdot \sqrt{\eps}$. On the hyperfinite Bernoulli experiment let

\[
Z=\frac{S-Np}{\sqrt{Nv}},\qquad
L_h=\left(\frac{p_h}{p}\right)^S
    \left(\frac{1-p_h}{1-p}\right)^{N-S},\qquad
\ell_h=\log L_h.
\]

Solve the following increasingly strong tasks.

\begin{enumerate}[leftmargin=*]
\item For fixed standard $h$ and finite $Z$, find $\st(\ell_h)$.
\item Fix standard $H>0$, choose unlimited $R$ with
   $1 \ll  R \ll  \sqrt{N}$, and prove uniformly for $|h|\le H$, $|Z|\le R$ that

   \[
   \ell_h=\frac{hZ}{\sigma}-\frac{h^2}{2v}
   +\sqrt\epsilon(1-2p)
      \left(\frac{h^3}{3v^2}-\frac{h^2Z}{2v^{3/2}}\right)
   +\rho_h,
   \quad |\rho_h|\le C_{p,H}\epsilon(1+|Z|).
   \]

\item Prove the good event $G_R=\{|Z|\le R\}$ has probability infinitesimally
   close to one uniformly on $|h|\le H$, prove $E_0[L_h]=1$, and establish
   mutual contiguity of the null and each fixed local alternative.
\item Prove the algebraic CLT, derive Le Cam's third-lemma shift
   $Z \Rightarrow  N(h/\sigma,1)$ under $p_{h}$, and find the asymptotic power of the
   one-sided level-$\alpha$ test.
\item Optional refinement: obtain the first continuity-corrected Bernoulli
   Edgeworth term and explain why omitting the half-cell correction leaves a
   lattice term of order $\sqrt{\eps}$.
\end{enumerate}

\textbf{Why this framework.} The null and local alternative have the same ordinary
parameter, but $\omega$ observations turn their $\sqrt{\eps}$ separation
into finite evidence. The first coefficient cancellation is short; the real
difficulty is constructing the entire experiment, proving a uniform analytic
remainder, changing probability algebraically, and retaining the first
finite-grid correction beyond Gaussian power.


\section*{Where the solutions are}

Worked solutions, with the corresponding checked Lean~4 kernels, are in the
companion booklet \href{https://files.pannous.com/hyperreal-exercise-solutions.pdf}{\emph{Twenty Exercises in Algebraic Probability ---
Solutions}}:
\begin{center}\small\url{https://files.pannous.com/hyperreal-exercise-solutions.pdf}\end{center}
Try the exercises first: the algebra is short, and the point of each one is
what survives that a real-valued shadow would have deleted.


\end{document}
